% =========================================================================
% Metaheuristics Project (31-MM34, SoSe 26) -- Report Structure
% Compile: pdflatex Report_Structure_V0.tex
% =========================================================================
\documentclass[10pt, a4paper]{article}

\usepackage[utf8]{inputenc}
\usepackage[T1]{fontenc}
\usepackage{lmodern}
\usepackage[margin=2.5cm]{geometry}
\usepackage{amsmath,amssymb}
\usepackage{booktabs}
\usepackage[hidelinks]{hyperref}

\title{An Adaptive Large Neighborhood Search for the\\
       Prize-Collecting Skill Vehicle Routing Problem}
\author{Umais Chaudhary, Rehman Kerimov, Dustin Schulz, Marco Leander vom Orde \\ \small 31-MM34 Metaheuristics, SoSe 26, Bielefeld University}
\date{Juli 2026}

\begin{document}
\maketitle

\begin{abstract}
% ~120 words: problem, chosen method (ALNS), key engineering insight
% (evaluation efficiency dominates in Python at 30s), headline result
% (63.0% of singleton UB as 3-seed mean, all feasible, on par with a
% heavily tuned single-trajectory LNS baseline).
\end{abstract}

% -------------------------------------------------------------------------
\section{Introduction}
% Motivation (pharmacy home delivery), problem features: optional customers,
% skills, time windows, shifts; 30s CPython budget.
% Red thread: literature -> ALNS is the right family; the actual battle is
% move-evaluation throughput; design principle "efficiency first, then
% operator richness". Contributions bullet list.

% -------------------------------------------------------------------------
\section{Problem Statement}
% Formal definition: sets C, B; profits, service times, time windows,
% shifts, skill sets; objective (1); feasibility rules; rounded Euclidean
% travel times. Singleton upper bound UB used as evaluation yardstick.

% -------------------------------------------------------------------------
\section{Related Work}
\subsection{Problem classification}
% TOPTW at the routing core; Skill VRP / technician routing for the skill
% aspect; NP-hardness.
\subsection{Metaheuristics for selective routing with time windows}
% ILS (Vansteenwegen), granular VNS (Labadie), I3CH (Hu & Lim), tuned
% ILS/SA (Gunawan); forward time slack as shared efficiency idea.
\subsection{LNS, ALNS and population methods}
% Shaw; Ropke & Pisinger; SISR; ALNS for TRSP (Kovacs); HGS (Vidal) and why
% it fits poorly here (selection + hard feasibility).
\subsection{Synthesis}
% Why ALNS with skill-aware insertion and SA acceptance is the evidence-
% backed choice under a short Python time budget.

% -------------------------------------------------------------------------
\section{Solution Approach}
\subsection{Framework overview}
% Pipeline + Algorithm 1 as package-free tabbing figure
% (ALNS main loop, time-based cooling, polish hook).
\subsection{Preprocessing and construction heuristic}
% Distance matrix, per-customer feasible-courier lists (skills + singleton
% time feasibility), profit-density bound; multi-start greedy insertion
% with four orderings.
\subsection{Destroy operators}
% Random, related (Shaw), worst-detour, worst-density, skill-scarcity;
% destroy size as part of operator identity (small/large variants).
\subsection{Repair operators}
% Greedy best-first, regret-2, sequential cheapest insertion
% (profit/random order), noised variant; candidate sampling
% (removed + top-density + random extras).
\subsection{Efficient move evaluation}
% The key enabler: (i) forward time slack -> O(1) feasibility,
% (ii) per-(customer,vehicle) move cache with per-route invalidation,
% (iii) early-exit suffix evaluation as fallback. Differential testing.
\subsection{Acceptance, adaptivity and intensification}
% SA with time-based cooling; T0 scaled to move deltas; adaptive weights
% normalized by operator runtime; return-to-best; randomized restarts.
\subsection{Local-search polish}
% Or-opt segment relocation + greedy refill of freed shift capacity on new
% global bests.

% -------------------------------------------------------------------------
\section{Experimental Investigation of the Components}
\subsection{Setup and evaluation metric}
% Hardware, CPython 3.9, 30s budget, seeds; %UB metric definition.
\subsection{Move-evaluation throughput}
% Ablation table on n1000: naive vs. +cache/early-exit vs. +slack fast path;
% quality follows throughput.
\subsection{Destroy size}
% Fixed vs. randomized fraction; size-split operator variants (table).
\subsection{Operator portfolio}
% Effect of related/worst-detour removal; effect of sequential repair;
% adaptive usage shares; full candidate pool tested and rejected.
\subsection{Acceptance design}
% T0 scale and time-based schedule; sensitivity remarks.
\subsection{Local-search polish and noised insertion}
% Marginal gains of the final two components.

% -------------------------------------------------------------------------
\section{Overall Performance}
\subsection{Final results}
% 20-instance table: UB, greedy start, ALNS 3-seed mean, %UB, baseline LNS.
\subsection{Comparison with a tuned single-trajectory LNS}
% Baseline description; expected-quality vs. lucky-run argument; 12/20
% wins/ties; n1000 dominated on every seed.
\subsection{Development progression and scaling}
% 61.0 -> 63.0 progression table; iteration counts vs. n; hardness driven
% by skill scarcity, not size; small-time-limit robustness.

% -------------------------------------------------------------------------
\section{Conclusion}
% Summary; main lesson (throughput-first engineering enables the adaptive
% framework); future work: SISR string removals, time-aware insertion
% scoring, set-covering post-optimization.

% -------------------------------------------------------------------------
\begin{thebibliography}{99}
% ~16 entries: Shaw 1998; Ropke & Pisinger 2006; Pisinger & Ropke 2007;
% Savelsbergh 1992; Vansteenwegen et al. 2009/2011; Labadie et al. 2012;
% Hu & Lim 2014; Gunawan et al. 2016/2017; Cappanera et al. 2011;
% Kovacs et al. 2012; Pillac et al. 2013; Cisse et al. 2017;
% Christiaens & Vanden Berghe 2020; Vidal et al. 2012/2022.
\end{thebibliography}

\end{document}
